\section{Third parties, and what the agents use talk for}
\label{sec:talk}

A trade between two players can change what a third player gets. Because a bankrupt player's deeds pass to its creditor, keeping a weak player solvent can keep a monopoly away from someone else, and letting it fail can deliver one. Of the \fact{d.eliminations.nonterminal} eliminations that left at least two players in the game, \fact{d.eliminations.nonterminal-to-eventual-winner} sent the estate to the eventual winner. The agents bargain over this prospect in so many words.

\IncidentTitle{g13}{A deed sold as insurance against a third player's inheritance}
At G13 turn 107 Grok has \$35, owns the yellow group and stands nine squares before Sol's Short Line, where it would owe \$200. It offers its mortgaged Pennsylvania Avenue to Opus for \$200:
\TranscriptQuote{g13-d250-insurance}
Opus counters at \$130. Grok holds at \$200 and says publicly that \$130 leaves it \$35 short. Its private note records that \$130 would be enough if it also mortgaged Ventnor. The provider-returned text of the same response states the omission:
\TranscriptQuote{g13-d252-omission}
Opus accepts at \$200 and pays \$16 of transfer interest. Its returned explanation takes up the estate argument:
\TranscriptQuote{g13-d253-buyer}
The inheritance the seller describes follows from the rules. Its public odds of 5 in 36 overstate a direct roll of nine, which is 4 in 36. Its own note uses the correct figure: \InlineQuote{in-g13-d252-floor}. With \$130 from the sale, \$35 in hand and the unmentioned \$130 mortgage, it could have paid the rent. Opus's text repeats the estate concern and does not address the financing claim. The records do not show that the omission produced the extra \$70. Grok missed Short Line on its next roll. Opus completed green at auction seven turns later. At turn 150 Grok did go bankrupt on Short Line, and the yellows passed to Sol, the outcome the sale had been pitched to prevent. Opus won.
\IncidentSource{G13 d250 to d253, with later events. Opus responds through action-only recovery, so its text is a provider-returned explanation. Grok's omission statement is also provider-returned text. Both of Grok's public messages appear in Opus's next prompts.}

The estate argument concerns a rival's options two steps removed: what Sol would gain if Grok failed. Both parties state it, the purchase agrees with it and its premise is correct. Like the refusals in Section~\ref{sec:rival}, it is level three.

Four further incidents, documented the same way in Appendix~\ref{app:cases}, show what else the agents do with public messages under a prompt that permits deception.

\emph{A bluff with a recorded ceiling.} In G7 Grok bids \$600 for a deed and calls the offer final, while its note says to accept \$625. The seller accepts \$600. Later the same seller, short of cash after building, sells a completing deed to Gemini for twice Gemini's bid. Gemini's note times the bid to the seller's need for building cash, the seller's note sets a price that funds its own houses and leaves Gemini little, and both build within eight turns.

\emph{A false premise corrected.} In G19 a Grok proposal reaches Sol with a claim that a pink deed is orange. Sol replies that the deed is pink and accepts the trade on its terms. Grok's own note in that response says to abort the trade, and the text outside its tool call names a different action.

\emph{A call to withhold cash.} In G20 Gemini asks the table not to buy a cash-poor rival's spare deeds, because the money would become houses on its new monopoly. Gemini then refuses such a sale itself. Opus later makes the opposite offer and describes it as money for more houses. No sale settles.

\emph{An assurance withdrawn.} In G18, with two players left, Opus tells Sol that its \$250 offer for a deed still stands. Text returned with that response says Opus had not raised the offer because extra cash would help the opponent last longer. One turn and one \$750 rent later, Sol tries to take the offer and Opus withdraws it, with the same stated reason. No Opus text in this stretch mentions inheriting anything. The stated motive is to deny the last opponent the cash to survive.

Each of these is one trajectory, chosen because it is clear. They are not a sample and do not measure how often agents bluff, omit, correct or coordinate.
